\chapter[Sermon 85. The Description of Christ the Judge Coming to the Last Judgment.]{The Description of Christ the Judge Coming to the Last Judgment.}
\label{sermon:085}
\markright{Sermon 85. The Description of Christ the Judge Coming to the Last ...}

And I saw heaven open and behold a white horse, and he that sat upon him was called Faithful and True, and in righteousness he judges and makes war. His eyes were as a flame of fire, and on his head were many crowns. And he had a name written that no man knew but himself. And he was clothed with a vesture dipped in blood, and his name is called the Word of God. And the warriors which were in heaven followed him upon white horses, clothed with white and pure silk. And out of his mouth went a sharp sword, that with it he should smite the heathen. And he shall rule them with a rod of iron, and he shall tread the winepress of the fierceness and the wrath of almighty God. And hath on his vesture, and on his thigh a name written: King of Kings, and Lord of Lords.

\index{Judgment, Last}Hitherto we have heard many things of the various punishments of the ungodly. And because it is manifest that God inflicts punishment on the mischievous and wicked at sundry times and in diverse ways, but most fully and most severely in that very last judgment, and from then on forever, and John has once, twice, thrice begun to treat of the last judgment, as in the end of chapters 11 and 14. And yet he has always deferred, suspended, and reserved it for another place; at last thinking it time to set before all men's eyes a description chiefly necessary, he takes it in hand and now finishes it as a matter of the greatest importance. He annexes therefore to a plentiful treatise of the torments of the ungodly, a most full and evident description of the judge most righteous and greatest, and of that last judgment, most strict of all others, wherein most fully and severely the pains shall be executed upon all Antichristians and ungodly forevermore. This place (which is the fourth of this chapter) and this treatise stretches unto chapter 21. The language is grand, smelling of prophetic majesty, apostolic clarity, and efficacy. You shall find not a few of this sort in the prophets, especially in chapters 24, 25, 26, and 27 of Isaiah.

And truly this doctrine is very profitable and necessary to be learned and understood most diligently by all faithful people. It was set forth with great diligence and abundant clarity to this end by the Prophets and Apostles, but especially by the Lord Jesus Christ himself, both in the Gospel and also in this most godly revelation. For unless you are kept in your duty by fear of the coming judgment and Judge, it is no surprise that you run mad and perish with this foolish and wicked world. In the treatise of the last judgment, one sees the end of all people, life and death, happiness and misery, pain and unspeakable, weighty reward. Whoever remembers these things well, abhors wickedness and walks in holy fear before God.

And we have learned of the doctrine of the Gospel, that the same day of the restoration of all, and oppression of the ungodly, and also of all ungodliness, is known to no mortal man, but to the Father alone: and therefore to inquire of the hour and moment thereof is most foolishly done, much more wickedly. Notwithstanding, the good Lord has shown and signified tokens, which when we shall see to be fulfilled and accomplished, we might lift up our heads, knowing that our redemption draws near. Behold your redemption, he says, not your torment. For he speaks of the godly, looking for their redemption from heaven, at the return of our Savior and Redeemer, our Lord Christ: which shall also inflict vengeance on his enemies, as Paul says in 2 Thessalonians 1. Therefore, let us not be curious here, searching for things unsearchable: but rather let us watch and pray, according to the wholesome precept of our Savior, judge, and revenger. Let us have our loins girded, and let lights burn in our hands. Let us look for him steadfastly in faith, and with holy hope. Let us rather take heed that the care of this world does not possess our hearts, and beware of drunkenness and surfeiting, and that we are not of the number or conversation of those who, in the days of Noah and Lot, regarded worldly things only, despised heavenly things, and laughed to scorn those who gave them good counsel, until the wrath of God was kindled and fell upon them when they least looked for it. We see all the tokens that are said should come before the day of our Lord to be fulfilled. Let us watch therefore: and these things considered in this wise, let us see and hear with great and diligent attention, what manner of judge of all shall come, and what that judgment shall be—most wished for by the godly, horrible and with trembling to be feared by the ungodly.

\index{Daniel (Book of)}First, John in the vision sees heaven open. For by a vision, to the end all things might be more evident, he not only tells so great a matter, but sets it before their eyes to behold: and that he says, he says of the revelation of Jesus Christ, lest any should object and say, “Are you not a mad fellow to talk thus of matters unknown?” For what is he that knows who or what that judge shall be? Or else what that judgment shall be? Therefore he tells these things from the judge Christ himself, and by a heavenly revelation. For other places of the Scripture show that the Lord shall come in glory and majesty: therefore with a great and most shining brightness of light, with fire and exceeding great clearness. For so it is said in Matthew 24:25 and 25:1, in Daniel 7, and in 2 Thessalonians 1. Therefore, by the opening of heaven is signified that the whole world shall be lightened with glory and brightness, and that the same day shall be most shining and clear. Others understand that the judgment cannot be fully perceived but through celestial revelation. Which as I confess to be most true, so I think I hear some greater matter to be signified.

Now follows the description of the judge, as of a noble and valiant warrior, composed of many parts. The faithful understand by this that the guardian, watchman, and avenger of the church does not sleep, whom the wicked do not perceive, failing to recognize the wrongs they commit against the godly, nor caring for those superstitious Christians, as they term them. They see moreover that they err if they think Christ at any time is too lenient, or overlooks for too long the calamities of his servants. For now he comes forth as a judge and revenger. There are many excellent descriptions of Christ in this book, as in any other; but this one is most elegant and vivid, which I have, according to my limited ability, expounded by parts. You shall always consider matters of greater import, until it is given to behold them directly with our eyes.

Our judge comes on horseback, and that on a white horse: not that he needs the help of corruptible horses in heaven, but thus he speaks after the manner of men, that we might imagine greater things. Conquerors ride on white horses. Here is signified therefore, that our judge shall be a conqueror and a triumpher. Others suppose the white horse to be a sign of his most pure humanity. I understand rather the white cloud. For the same token he took up from the eyes of his disciples, when he ascended into heaven from Mount Olivet. In the same way he shall come again to judge. And just as kings are carried on horses and chariots, so the Psalmist ascribes to God clouds as horses and chariots.

Our judge is faithful and true. Faithful to his faithful. True in all his promises toward the godly and ungodly. They are deceived, and shall see themselves to be deceived at the judgment, so many as have contemned the promises and threatenings of God as vain, and esteeming things after the success of this world, judge the wicked to be happy and fortunate, and the godly to be wretched and miserable. Hereof hath the Prophet Malachi reasoned in the third and fourth chapters. And saying the judge is faithful and true, he judges and fights in righteousness, to wit, giving every man his own, rewards to the good, and punishments to the evil. This king doth not judge and fight as the kings of this world are wont, following vanity and corrupt affections. And Christ is said to fight when he rewards the ungodly after their demerits, the Apostle: we must all, says he, be manifest before the judgment seat of Christ, that every man may receive such things as he hath done by his body, according to that he hath wrought be it good or evil. 2 Corinthians 5.

3. The eyes of the Judge are like a flame of fire. For just as no one can escape or hide from a judge or judgment, as he searches the secrets of all, and nothing can be hidden from his sight, so are his eyes terrible and fearful against the ungodly. The godly, however, are filled with all pleasure, joy, and gladness by the sight of the Lord. Flaming and fiery eyes are also attributed to Christ in the first vision; see there for more. And Scripture testifies everywhere that the judge knows all things, even the secrets of hearts. Therefore, you act foolishly if you think you have won the battle and have sinned unpunished, having escaped the knowledge and judgment of men. Another judgment remains, in which all the deeds of the wicked will be revealed to their utter shame and confusion before all the world. The sins of the godly are covered by him through whose benefit they are justified and absolved from pain and guilt.

\index{Rome}4. Our Judge has very many crowns upon his head, for he alone governs all realms and nations. As Daniel signified in chapter 7. He alone might truly be called Africanus, Europeanus, and Asiaticus, Parthicus, Persicus, Germanicus, Gothicus, and others. Which our kings have foolishly claimed for themselves, attempting to assume the monarchy, when Christ alone is the true Monarch forever. This Judge and mighty Prince shall strike off the triple crown from the head of the Bishop of Rome. Moreover, there shall be no king so mighty in the whole world that shall be able to resist him, and make war against him.

5. Our judge has a name written, which no man knows save himself. This will be explained more clearly in due course. Christ has a name unspeakable, for he is the true God, eternal, incomprehensible, and Almighty. This name knows no one but himself. For first, the majesty of God is greater than can be comprehended by any creature; again, the name of God is agreeable to no one but him alone, for the name of God, in this signification, may not be communicated. For he is very God, and besides him none, which thing Isaiah repeats often. He is the Savior, King, Monarch, and Judge: these things all belong properly to him alone and are not common to others. Moreover, the Lord himself says in the gospel: no man has known the Son but the Father; neither has any man known the Father, save the Son, and to whom the Son has pleased to reveal. Besides this, we see it imperfectly: and the glory of the divine majesty is so great, as I said before, that human capacity is unable to conceive such a glory. No man therefore save God alone knows his name.

6. The vesture of our Judge was sprinkled with blood. This signifies victory and the slaughter of his enemies, which will be explained further toward the end of the chapter. He took this observation about our Judge from chapter 63 of \textit{Explanations}. He alludes to conquerors returning from battle, whose garments and armor are soaked with the blood of the slain. And it betokens the just severity of the Judge and great slaughter of the enemies.

7. The name of the judge is now expressed, [?which] is utterly unknown to the ungodly. And the Judge is called the word of God. For the Son is the word and speech of God, the express mark of the divine substance: in whom the Father himself is expressed: and of whom, as of the word, we understand the will and mind of the Father, the true messenger of the heart. These holy words of the gospel are known: In the beginning was the word, [?and] the word was with God, \&c. Therefore Christ, the word, was made flesh, the Lord God and Judge of all.

\index{John, Saint}8. To the Judge is added an Army, not of angels only, with whom he often repeated in the gospel that he would come in judgment: but of all the faithful, or saints, who at no time, not even here, are separated from their head. For at the sound of the trumpet, the archangel arises, and the saints arise, and the living along with the dead are changed and are taken up to meet Christ in the air. Here, here in the clouds and bright air, appear with Christ the happy and blessed victors. Immediately, the ungodly also rise, and those who lived at that time are changed with those who rise again, to pain and confusion. But they see the saints with Christ in heaven, in glory, and immediately feel unspeakable torments. These things, without doubt, come to pass and are fulfilled, as described in the 3rd and 5th chapters of Wisdom. Saint John therefore says that this army is in heaven, not on earth. He says how they follow Christ. For the same thing the Apostle also said in the first letter to the Thessalonians, the 4th. Moreover, he adds that they were clothed, and did not appear naked, and he expresses the kind of garment. They were clothed (he says) in silk, white and clean. For saints in Christ obtain righteousness and glory, are made clean and are glorified. And this meaning Saint John himself has opened to us a little before, saying: silk is the righteousness of saints.

9. From the Judge’s mouth proceeds a two-edged and sharp sword, which cuts on either side. It is not sharp on one side and blunt on the other; it cuts on both sides equally. This signifies a just sentence pronounced by God’s mouth against the wicked. For against them, God’s sentence is a sword, piercing even to their hearts. Therefore it is also called sharp. The judgment of our Judge is straight and severe, but yet just and righteous. What that sword is, is declared in the Gospel: namely, that heavy and immutable sentence, “Get you hence into everlasting fire.” Matthew 25. Upon this it follows in the words of the Evangelist: that with the same he may strike the heathen, that is, to condemn and put all unbelievers to perpetual torments.

\index{Repentance}10. And he shall rule them with a rod of iron. In the same way, he says the same thing he said before. For those who would not receive or acknowledge with repentance the staff of instruction and discipline afterward, shall find judgment and feel the iron scepter, with which he will break them all to pieces, like a potter’s vessel. Neither shall any power resist or prevail against him. And this manner of speaking is taken from Psalm 110. For Saint John gladly uses the words of Scripture to the end to make his book more commendable or more pleasant and acceptable.

11. He treads the winepress of the wine of wrath. And again he says the same, that he did before but now utters another parable, taken from the scriptures, namely from the sixty-third chapter of Isaiah. The effect of it, or some part of it, is that he will pour out his wrath upon the ungodly and punish them most severely with his almighty hand, to which all things give way, bowing their heads. See what is said here in the fourteenth chapter of this book.

12. Again, the name of this judge is shown, and in the name is the majesty and power of all others greatest. He has the name written on his garment and on his thigh. By these is declared the true humanity of Christ, after which he is exalted, as the Apostle says in the second chapter of Philippians. And to him is given a name which is above all names. Here he is called King of Kings and Lord of Lords, very God, Lord, monarch, and judge of all men. For so do the other Apostles speak also in the second chapter and seventeenth verse of Acts. And there might seem in this name of the Judge, as it were a cause to be shown, wherefore he is here appointed judge over all, because he is King and Lord of all. To whom be glory forever. Amen.
